\chapter{Interventional Effects}

Motivated by the problems presented by situations with treatment-induced confounding,~\cite{vanderweeleetal2014} 
propose a randomized-interventional analogue of natural effects that are identified without 
requiring the cross-world independence assumption. Instead of defining 
a unit-level direct and indirect effect estimand as in the case of natural effects, interventional effects
target population-level average outcomes from a hypothetical joint intervention on the treatment and the 
mediator. As mentioned in Section~\ref{chapter:sample}, a key difference between natural and interventional effects
is that interventional effects can still be identified given a set of reasonable assumptions in cases where
there is either treatment-induced confounding or when there are multiple mediators of interest. We begin by
reviewing the definition and interpretation of interventional effects in one mediator problems before 
extending our analysis to the multiple mediator case. We conclude this section with a discussion of what we
view as the practical use and interpretation of interventional effects. 

\section{One Mediator}

\subsection{Effect Definitions}\label{subsection:intv_eff_defs}

Let $G_{i}(a)$ represent a random draw of mediator value from $M(a)$, which now represents a 
the population-level distribution of the potential mediator given treatment level $a$ (and covariates $C_{i}$).
Let $Y_{i}(a,G_{i}(a'))$ represent a potential outcome for unit $i$ assigned treatment $a$ and
a random level for the mediator from the distribution of $G(a')$. This is a stochastic potential 
outcome: the potential outcome is itself a random variable because its value depends on a realization
from the distribution of $G(a')$~\citep{vanderweelerobins2012}. 

The definitions for interventional effects mirror those for natural effects above, except that they are
defined in terms of stochastic potential outcomes $Y_{i}(a,G_{i}(a'))$ that represent joint interventions
to assign unit $i$ values for the treatment and the mediator. First, we write out the equation to estimate
the mean of $Y_{i}(a,G_{i}(a'))$ for any $a, a' \in \mathcal{A}$~\citep{kuhaetal2021}: 
% Should try to mirror this definition in the natural effects section

\begin{equation}\label{eq:RIE_one_mediator_g_comp}
\begin{aligned}
	\mathbb{E}[Y_{i}(a,G_{i}(a'))|C_{i}] = 
	\int_{m} \mathbb{E}[Y_{i}(a,m)|A_{i} = a, M_{i} = m, C_{i}]\mathbb{P}[M_{i}(a') = m|C_{i}]dm
\end{aligned}
\end{equation}

We then define interventional effects in the following way: 

\begin{itemize}[noitemsep,topsep=0pt]
\item  $\mathbb{E}[Y_{i}(1,G_{i}(1)) - Y_{i}(0,G_{i}(0))]$ --- The \textit{interventional total effect} 
represents the effect of a joint intervention that gives a unit treatment while assigning their mediator 
level to a random draw from the distribution of the mediator for treated units, compared to an intervention
to assign a unit to no treatment while drawing their mediator level from the distribution of the mediator for 
untreated units. 
\item  $\mathbb{E}[Y_{i}(1,G_{i}(0)) - Y_{i}(0,G_{i}(0))]$ --- The \textit{pure interventional direct effect} 
represents the effect of a joint intervention that gives a unit treatment while drawing their mediator level 
from the distribution of the mediator for untreated units, compared to an intervention
to assign a unit to no treatment while drawing their mediator level from the distribution of the mediator for 
untreated units. 
\item  $\mathbb{E}[Y_{i}(1,G_{i}(1)) - Y_{i}(0,G_{i}(1))]$ --- The \textit{total interventional direct effect} 
represents the effect of a joint intervention that gives a unit treatment while drawing their mediator level 
from the distribution of the mediator for treated units, compared to an intervention
to assign a unit to no treatment while drawing their mediator level from the distribution of the mediator for 
treated units. 
\item  $\mathbb{E}[Y_{i}(0,G_{i}(1)) - Y_{i}(0,G_{i}(0))]$ --- The \textit{pure interventional indirect effect} 
represents the effect of a joint intervention that gives a unit no treatment while drawing their mediator level 
from the distribution of the mediator for treated units, compared to an intervention
to assign a unit to no treatment while drawing their mediator level from the distribution of the mediator for 
treated units.
\item  $\mathbb{E}[Y_{i}(1,G_{i}(1)) - Y_{i}(1,G_{i}(0))]$ --- The \textit{total interventional indirect effect} 
represents the effect of a joint intervention that gives a unit treatment while drawing their mediator level 
from the distribution of the mediator for treated units, compared to an intervention
to assign a unit to treatment while drawing their mediator level from the distribution of the mediator for 
treated units.
\end{itemize}

The above effects are identified under consistency, composition, and conditions (i)-(iii) of sequential ignorability. 
In particular, the cross-world independence assumption~\citep{vanderweeleetal2014} is not required. 
The interventional total effect and the total effect are equivalent when the cross-world independence
assumption is satisfied, but the two quantities are not otherwise equivalent~\citep{vanderweeleetal2014}.
Conceptually, this occurs because the joint intervention on the treatment and the mediator breaks the 
causal "link" between $A$ and $M$ that is preserved when estimating the total causal effect.

% Might need to mention Vanderweele-Robinson interventional effects here: these seem to cover the pure and total 
% interventional effects (and allow us to not have assumptions (i) and (ii)); the effect definitions would be 
% identical (though of course only the indirect effects would exist) and so the only thing that changes is the 
% interpretation; another possibility is to put this in the "interpretation of effects" section, which has the 
% advantage of not cluttering the technical introduction, and I would not lose much in terms of interpreting 
% these effects

Much like with natural effects, there is a decomposition of the interventional total effect into the sum
of the pure interventional direct effect and the total interventional indirect effect: 

\begin{equation}
\begin{aligned}
	\mathbb{E}[Y_{i}(1,G_{i}(1)] - \mathbb{E}[Y_{i}(0,G_{i}(0))] = 
	\{\mathbb{E}[Y_{i}(1,G_{i}(0))] - \mathbb{E}[Y_{i}(0,G_{i}(0))]\}\\ 
	+ \{\mathbb{E}[Y_{i}(1,G_{i}(1))] - \mathbb{E}[Y_{i}(1,G_{i}(0))]\}
\end{aligned}
\end{equation}

Following the decomposition for natural effects, the most common way of writing this 
decomposition is in terms of the pure direct and total
indirect effects, but it is again easy to see that a decomposition also exists in terms of the total direct
and pure indirect effects. The pure and total interventional effects are not explicitly defined 
by~\cite{vanderweeleetal2014} and a distinction is generally not made between them in the literature, though 
the notion of a pure and total interventional effect appears in the eAppendix of~\cite{vanderweele2014}. 
Additionally,~\cite{didelez2006} and ~\cite{geneletti2007} defined a form of these types of
interventional effects in terms of general randomized interventions on the mediator without defining these
effects in terms of counterfactuals. These allow for any stochastic intervention on the mediator rather than
restricting such interventions to the distribution of the mediator in two treatment groups. The 
total interventional effect is formally defined in~\cite{linvanderweele2017}. 

~\cite{vanderweele2014} also gives a 3-way decomposition of the total interventional effect in terms of the 
pure interventional direct and indirect effects and a third term that can be interpreted as the causal effect
of an interaction between the treatment and the mediator:

\begin{equation}
\begin{aligned}
	\mathbb{E}[Y_{i}(1,G_{i}(1))] - \mathbb{E}[Y_{i}(0,G_{i}(0))] = 
	\{\mathbb{E}[Y_{i}(1,G_{i}(0))] - \mathbb{E}[Y_{i}(0,G_{i}(0))]\}\\ 
	+ \{\mathbb{E}[Y_{i}(1,G_{i}(1))] - \mathbb{E}[Y_{i}(1,G_{i}(0))]\}\\
	= \{\mathbb{E}[Y_{i}(1,G_{i}(0))] - \mathbb{E}[Y_{i}(0,G_{i}(0))]\}
	+ \{\mathbb{E}[Y_{i}(0,G_{i}(1))] - \mathbb{E}[Y_{i}(0,G_{i}(0))]\} + \\
	\{(\mathbb{E}[Y_{i}(1,G_{i}(1))] - \mathbb{E}[Y_{i}(1,G_{i}(0))]) - 
	(\mathbb{E}[Y_{i}(0,G_{i}(1))] - \mathbb{E}[Y_{i}(0,G_{i}(0))])\}
\end{aligned}
\end{equation}

% -------
% This formula only makes sense if there is a binary mediator (such that m* = 0, m = 1); it is also generally
% unclear and does not add very much to my discussion here; consider cutting; also, what actually is the 
% interpretation of the interaction term here and how does it differ from the interaction term in the natural
% effects decomposition? Is the decomposition really the right tool here if I want to emphasize that these
% decompositions are not particularly important? Or should I just offer them as a helpful analytical tool to 
% help with a taxonomy of multiple mediator interventional effects?
% (this also will motivate their inclusion in the natural effects section)
Writing out the third term of this decomposition, we get from the eAppendix of~\cite{vanderweele2014} 

\begin{equation}
\begin{aligned}
	\{(\mathbb{E}[Y_{i}(1,G_{i}(1))] - \mathbb{E}[Y_{i}(1,G_{i}(0))]) - 
	(\mathbb{E}[Y_{i}(0,G_{i}(1))] - \mathbb{E}[Y_{i}(0,G_{i}(0))])\} =\\
	\int_{m \in \mathcal{M}} \mathbb{E}[(Y(a,m) - Y(a^{*},m)) - (Y(a,m^{*}) - Y(a^{*},m^{*}))|c]*\\
	(\mathbb{P}[M(a) = m | c] - \mathbb{P}[M(a^{*}) = m | c)])
\end{aligned}
\end{equation}
% -------

The total interventional effect and the total causal effect are equivalent when the cross-world independence
assumption is satisfied, but the two quantities are not otherwise equivalent~\citep{vanderweeleetal2014}.
Conceptually, this occurs because the joint intervention on the treatment and the mediator breaks the 
causal "link" between $A$ and $M$ that is preserved when estimating the total causal effect.
% Derivation of this result

\subsection{Experimental Identification}

A key advantage of interventional effects compared to natural effects is that it is in principle possible to 
design randomized experiments that can identify these effects. Although such experiments might be infeasible 
or unethical in practice, being able to conceptualize an ideal experiment to identify such effects can still
help inform observational study design and interpretation in many ways. The ways in which experimental 
identification aids in the interpretation of interventional effects is expounded upon in 
Section~\ref{subsection:interpretation}. 

The experimental setup, like our running example, is inspired by~\cite{morenobetancur2018}. Suppose we 
are trying to estimate interventional effects in the running example. Suppose a "headache drug" has been developed 
that allows us exactly set the level of headache a patient who takes this drug will experience while having no 
other side effects. We set up the following 6-arm trial, with patients randomly assigned to each arm: 

\begin{experiment}\label{exp:one_mediator}
\indent 
\begin{enumerate}[noitemsep,topsep=0pt]
	\item Patients in \textbf{Arm 1} receive the anti-depressant, while patients in \textbf{Arm 2} do not. 
	The levels of both headache and depression are recorded for each patient.
	\item After recording the results of the experiment in stage 1, patients in arms 3-6 receive the following
	combinations of treatment and headache medication: 
	\begin{enumerate}[noitemsep,topsep=0pt]
		\item \textbf{Arm 3} receives the anti-depressant, and the headache drug is used to give each patient 
		a randomly selected level of headache from the distribution of observed headache levels in Arm 1. 
		\item \textbf{Arm 4} receives the anti-depressant, and the headache drug is used to give each patient 
		a randomly selected level of headache from the distribution of observed headache levels in Arm 2. 
		\item \textbf{Arm 5} does not receive the anti-depressant, and the headache drug is used to give each patient 
		a randomly selected level of headache from the distribution of observed headache levels in Arm 1. 
		\item \textbf{Arm 6} does not receive the anti-depressant, and the headache drug is used to give each patient 
		a randomly selected level of headache from the distribution of observed headache levels in Arm 2. 
	\end{enumerate}
\end{enumerate}
\end{experiment}

% Table 1 -------
% Need to figure out dynamic labeling/referencing
\vspace{10pt}
\noindent
\begin{tabularx}{\textwidth}{@{}
  >{\hsize=0.40\hsize\centering\arraybackslash}X
  >{\hsize=0.40\hsize\centering\arraybackslash}X
  >{\hsize=1.40\hsize\centering\arraybackslash}X
  >{\hsize=1.40\hsize\centering\arraybackslash}X
  >{\hsize=1.40\hsize\centering\arraybackslash}X @{}}
\multicolumn{5}{@{}p{\textwidth}@{}}{%
  \textbf{TABLE 1.}\quad RCT detailed in Experiment~\ref{exp:one_mediator}} \\
\addlinespace[4pt]
\hline
\addlinespace[6pt]
\makecell[bc]{\textbf{Stage}} &
\makecell[bc]{\textbf{Arm}} &
\makecell[bc]{\textbf{Treatment}} &
\makecell[bc]{\textbf{Headaches}} &
\makecell[bc]{\textbf{Depression}} \\
\addlinespace[3pt]
\noalign{\hrule height 1pt}
\addlinespace[3pt]
\multirow{2}{*}{1} 
&1 & 1 & $\hat{M}(1)_{1}$ & $y_{1}$\\
&2 & 0 & $\hat{M}(0)_{2}$ & $y_{2}$\\
\addlinespace[2pt]
\noalign{\hrule height 0.5pt}
\addlinespace[2pt]
\multirow{4}{*}{2}
&3 & 1 & $\hat{G}(1)_{3}$ & $y_{3}$\\
&4 & 1 & $\hat{G}(0)_{4}$ & $y_{4}$\\
&5 & 0 & $\hat{G}(1)_{5}$ & $y_{5}$\\
&6 & 0 & $\hat{G}(0)_{6}$ & $y_{6}$\\
\addlinespace[3pt]
\noalign{\hrule height 1pt}
\end{tabularx}

\vspace{4pt}
{\footnotesize
\emergencystretch=3em
\setlength{\parindent}{1.5em}
\textsuperscript{1} $\hat{M}(1)$ and $\hat{M}(0)$ represent the distributions of headache in
arms 1 and 2, respectively. $\hat{G}(1)$  and $\hat{G}(0)$ represent the distributions of assigned
headache from random draws from $\hat{M}(1)$ and $\hat{M}(0)$, respectively. 

Finally, $y$ gives the observed average level of depression in each arm.\par}
\vspace{10pt}
% -------

Given this experimental setup, it is possible to identify the suite of interventional effects defined in 
Section~\ref{subsection:intv_eff_defs} by taking observed differences in outcomes $y$ (see Table 1). For
example, the estimated total interventional effect is $y_{3} - y_{6}$, the pure interventional direct
effect is $y_{4} - y_{6}$, and the pure interventional indirect effect is $y_{5} - y_{6}$. 

% I should consider introducing the Vanderweele-Robinson effects here, instead of first mentioning them
% in the multiple mediators (rename to "two mediators"?) section

\section{Multiple Mediators}

We now extend the interventional effects framework to a situation with two mediators. Additionally, we 
assume that we know the causal ordering of these two mediators (though this condition can be relaxed in 
some cases). Extending the running example from~\cite{morenobetancur2018}, we now suppose that we notice 
that the anti-depressant, in addition to causing headache, can cause insomnia in the first few weeks of
treatment. We also believe that insomnia can cause chronic headaches, so headaches can be caused both 
directly by the anti-depressant and indirectly through the anti-depressant's effect on insomnia. 
% Maybe should add a DAG here; it is probably much easier/more helpful to parse when presented right after
% the empirical example, even though it requires the assumptions for strict correctness; but I am using
% the DAGs for illustrative purposes anyway, and they can also represent the "causal story" we are interested
% in testing

In general, there are a number of different definitions of interventional effects with
two mediators~\citep{vanderweeleetal2014,vansteelandtdaniel2017,linvanderweele2017,jacksonvanderweele2018}
in the literature, each of which has meaningfully different identification assumptions and interpretations.
The main goal of this section is to introduce a unified framework for interventional effects that can 
accomodate each of these different effect definitions and provide a clear way to compare the assumptions 
and interpretations of these different effects. The motivation for this framework comes from the ability to 
devise a family of hypothetical experiments that can identify these effects. 

\subsection{A Unified Framework for Interventional Effects with Two Mediators}

We now let $M_{1,i}$ be the observed level of insomnia for unit $i$ and $M_{2,i}$ the observed level of
headache. Let $\mathcal{M}_{1}$ be the support of $M_{1,i}$ and $\mathcal{M}_{2}$ the support of $M_{2,i}$. 
We also redefine potential mediators and outcomes as follows:

\begin{itemize}[noitemsep,topsep=0pt]
	\item Let $M_{1,i}(a)$ represent the potential levels of insomnia for unit $i$ given 
	treatment level $a$. 
	\item Let $M_{2,i}(a,M_{1,i}(a'))$ represent the potential levels of headache for unit $i$ given 
	treatment level $a$ and an insomnia level $M_{1,i}(a')$ that they would have if they received 
	treatment level $a'$. 
	\item Let $Y_{i}(a,M_{1,i}(a'),M_{2,i}(a'',a'''))$ represent the potential outcomes 
	for unit $i$ given treatment level $a$ and potential mediator levels $M_{1,i}(a')$ and 
	$M_{2,i}(a'',a''')$. 
\end{itemize}

Finally, we define $G_{1,i}(a)$ as a random draw from the distribution of $M_{1}(a)$ 
and\\ $G_{2,i}(a,M(a'))$ as a random draw from the distribution of $M_{2}(a,M(a'))$. 

We conceive of interventional 
effects with two mediators to be the average potential outcome of a joint intervention that assigns
a treatment level $a$, a level of the fist mediator randomly drawn from the distribution of the 
mediator in a population with treatment level $a'$, and a level of the second mediator randomly drawn
from the distribution of the mediator in a population with treatment level $a''$ and first mediator 
level randomly drawn from a population with treatment level $a'''$. 
We now give a formula for the population mean of $Y_{i}(a,M_{1,i}(a'),M_{2,i}(a'',a'''))$ that
extends~\ref{eq:RIE_one_mediator_g_comp} to the two mediator case. This is similar to the formula given
by~\cite{linvanderweele2017}. 

\begin{equation}\label{eq:RIE_two_mediator_g_comp}
\begin{aligned}
	\Psi(a,a',a'',a''') = \mathbb{E}[Y_{i}(a,G_{i}(a'))|C_{i}] =\\ 
	\int_{m_{2}}\int_{m_{1}'}\int_{m_{1}} \mathbb{E}[Y_{i}(a,m_{1},m_{2})|A_{i} = a, 
	M_{1,i} = m_{1}, M_{2,i} = m_{2}, C_{i}]\\
	\mathbb{P}[M_{2,i}(a'',m') = m_{2}|A_{i} = a''', M_{i} = m_{1}', C_{i}]\\
	\mathbb{P}[M_{1,i}(a''') = m_{1}|C_{i}]\mathbb{P}[M_{1,i}(a') = m_{1}'|C_{i}]dm_{1}dm_{1}'dm_{2}
\end{aligned}
\end{equation}
% Should comment on the "two versions" of M1 here; can highlight this as the reason path-specific natural 
%effects are never identified

Following the one mediator case, we define interventional effects in terms of comparisons between different 
sequences of the interventions on treatment, first mediator, and second mediator we mentioned above. We
give a few examples of such effects below: 

\begin{itemize}[noitemsep,topsep=0pt]
	\item $\Psi(1,1,1,1) - \Psi(0,0,0,0)$ --- The \textit{interventional total effect} gives the 
	represents the effect of a joint intervention that gives a unit treatment while assigning $M_{1}$
	to a random draw from the distribution of $M_{1}(1)$ and $M_{2}$ to 
	a random draw from the distribution of $M_{2}(1,1)$, compared to an intervention
	to assign a unit to no treatment while drawing $M_{1}$ from $M_{1}(0)$ and $M_{2}$ from $M_{2}(0,0)$. We 
	call the first joint intervention the \textit{full treatment} regime and the second joint intervention the 
	\textit{baseline control} regime. 
	\item $\Psi(0,1,0,0) - \Psi(0,0,0,0)$ --- This gives the effect of changing the joint intervention from
	the baseline control regime described in the first bullet to an intervention that assigns $M_{1}$ 
	from the distribution of $M_{1}(1)$ but holds the other two interventions constant. 
	\item $\Psi(1,1,1,1) - \Psi(1,1,0,1)$ --- This gives the effect of changing the joint intervention to the
	the full treatment regime described in the first bullet from an intervention that assigns $M_{1}$ 
	from the distribution of $M_{1}(0)$ but holds the other two interventions constant at the full treatment
	levels. 
	\item $\Psi(0,0,0,1) - \Psi(0,0,0,0)$ --- 
\end{itemize}

There are many estimands defined in terms of comparisons of two values of $\Psi(a,a',a'',a''')$ that we might
be interested in, and which estimands we are interested in should
be driven by substantive considerations (see Section~\ref{subsection:interpretation}). For example, 
there is no \textit{a priori} reason why an interventional effect defined as a comparison between a
particular joint intervention and the baseline control regime is more meaningful than a comparison between
the same intervention and the full treatment regime, or indeed any other reference intervention. 

~\cite{linvanderweele2017} proved that equation~\ref{eq:RIE_two_mediator_g_comp} identifies 
$\Psi(a,a',a'',a''')$, and thus any estimands defined in terms of $\Psi(a,a',a'',a''')$, 
under the following expanded set of ignorability assumptions: 
% Also mention Vanderweele and Tchetgen Tchetgen 2017

\begin{assumption}\label{assm:seq_ignorability_2med} 
\indent
\begin{enumerate}[noitemsep,topsep=0pt,label=(\roman*)]
	\item $Y_{i}(a,m_{1},m_{2}) \indep A_{i} | C_{i}$ for all 
	$a \in \mathcal{A}$, $m_{1} \in \mathcal{M}_{1}$, $m_{2} \in \mathcal{M}_{2}$. 
	\item $M_{1,i}(a) \indep A_{i} | C_{i}$ for all $a \in \mathcal{A}$. 
	\item $M_{2,i}(a) \indep A_{i} | C_{i}$ for all $a \in \mathcal{A}$. 
	\item $Y_{i}(a,m_{1},m_{2}) \indep M_{2,i} | C_{i}$ 
	for all $a \in \mathcal{A}$, $m_{1} \in \mathcal{M}_{1}$, $m_{2} \in \mathcal{M}_{1}$. 
	\item $Y_{i}(a,m_{1},m_{2}) \indep M_{1,i} | C_{i}$ 
	for all $a \in \mathcal{A}$, $m_{1} \in \mathcal{M}_{1}$, $m_{2} \in \mathcal{M}_{1}$. 
	\item $M_{2,i}(a,m_{1}) \indep M_{1,i} | C_{i}$ 
	for all $a \in \mathcal{A}$, $m_{1} \in \mathcal{M}_{1}$
\end{enumerate}
where $0 < \mathbb{P}[A_{i} = a | C_{i}]$, $0 < \mathbb{P}[M_{1,i} = m_{1} | A_{i} = a, C_{i}]$,
$0 < \mathbb{P}[M_{2,i} = m_{2} | A_{i} = a, M_{i} = m_{1}, C_{i}]$
for all $a \in \mathcal{A}$, $m_{1} \in \mathcal{M}_{1}, m_{2} \in \mathcal{M}_{2}$. 
\end{assumption}

Assumption~\ref{assm:seq_ignorability_2med} is a simple generalization of parts (i)-(iii) of 
Assumption~\ref{assm:seq_ignorability} to the two causally ordered mediators case. It rules out
any treatment-mediator confounding and any mediator-mediator confounding after controlling for 
pre-treatment covariates. 

We now give the details of a hypothetical experiment, which is an extension of Experiment~\ref{exp:one_mediator} 
that can identify any estimands of interest that are defined in terms of $\Psi(a,a',a'',a''')$. 
We suppose that in addition to the "headache drug" defined for the previous experiment, there 
exists a similar "insomnia drug" that allows us to exactly set the level of insomnia a patient
experiences. We set up the following 22-arm trial (see Table 2):  

\begin{experiment}\label{exp:two_mediators}
\indent 
\begin{enumerate}[noitemsep,topsep=0pt]
	\item Patients in \textbf{Arm 1} receive the anti-depressant, while patients in \textbf{Arm 2} do not. 
	The levels of insomnia, headache, and depression are recorded for each patient.
	\item After recording the results of the experiment in stage 1, patients in \textbf{Arms 3-6} 
	receive combinations of treatment levels and are assigned a randomly selected level of insomnia 
	from the distribution of observed insomnia levels in either Arm 1 or Arm 2, similar to in 
	Experiment~\ref{exp:one_mediator}. The levels of insomnia, headache, and depression are recorded 
	for each patient.
	\item After recording the results of the experiment in stage 2, patients in \textbf{Arms 7-22} 
	receive combinations of treatment levels, a randomly selected level of insomnia from the distribution of
	observed insomnia levels in either Arm 1 or Arm 2, and a selected level of headache from the 
	distribution of observed headache levels in one of Arms 3-6. 
\end{enumerate}
\end{experiment}

% Table 1 -------
% Need to figure out dynamic labeling/referencing
\vspace{10pt}
\noindent
\begin{tabularx}{\textwidth}{@{}
  >{\hsize=0.50\hsize\centering\arraybackslash}X
  >{\hsize=0.50\hsize\centering\arraybackslash}X
  >{\hsize=1.25\hsize\centering\arraybackslash}X
  >{\hsize=1.25\hsize\centering\arraybackslash}X
  >{\hsize=1.25\hsize\centering\arraybackslash}X
  >{\hsize=1.25\hsize\centering\arraybackslash}X @{}}
\multicolumn{6}{@{}p{\textwidth}@{}}{%
  \textbf{TABLE 2.}\quad RCT detailed in Experiment~\ref{exp:two_mediators}} \\
\addlinespace[4pt]
\hline
\addlinespace[6pt]
\makecell[bc]{\textbf{Stage}} &
\makecell[bc]{\textbf{Arm}} &
\makecell[bc]{\textbf{Treatment}} &
\makecell[bc]{\textbf{Insomnia}} &
\makecell[bc]{\textbf{Headaches}} &
\makecell[bc]{\textbf{Depression}} \\
\addlinespace[3pt]
\noalign{\hrule height 1pt}
\addlinespace[3pt]
\multirow{2}{*}{1} 
&1 & 1 & $\hat{M_{1}}(1)_{1}$ & $\hat{M_{2}}(1,\hat{M_{1}}(1)_{1})_{1}$ & $y_{1}$\\
&2 & 0 & $\hat{M_{1}}(0)_{2}$ & $\hat{M_{2}}(0,\hat{M_{1}}(0)_{2})_{2}$ & $y_{2}$\\
\addlinespace[2pt]
\noalign{\hrule height 0.5pt}
\addlinespace[2pt]
\multirow{4}{*}{2}
&3 & 1 & $\hat{G}_{1}(1)_{3}$ & $\hat{M_{2}}(1,\hat{G}_{1}(1)_{3})_{3}$ & $y_{3}$\\
&4 & 1 & $\hat{G}_{1}(0)_{4}$ & $\hat{M_{2}}(1,\hat{G}_{1}(0)_{4})_{4}$ & $y_{4}$\\
&5 & 0 & $\hat{G}_{1}(1)_{5}$ & $\hat{M_{2}}(0,\hat{G}_{1}(1)_{5})_{5}$ & $y_{5}$\\
&6 & 0 & $\hat{G}_{1}(0)_{6}$ & $\hat{M_{2}}(0,\hat{G}_{1}(0)_{6})_{6}$ & $y_{6}$\\
\addlinespace[2pt]
\noalign{\hrule height 0.5pt}
\addlinespace[2pt]
\multirow{16}{*}{3}
&7 & 1 & $\hat{G}_{1}(1)_{7}$ & $\hat{G_{2}}(1,\hat{G}_{1}(1)_{3})_{7}$ & $y_{7}$\\
&8 & 1 & $\hat{G}_{1}(1)_{8}$ & $\hat{G_{2}}(1,\hat{G}_{1}(0)_{4})_{8}$ & $y_{8}$\\
&9 & 1 & $\hat{G}_{1}(1)_{9}$ & $\hat{G_{2}}(0,\hat{G}_{1}(1)_{5})_{9}$ & $y_{9}$\\
&10 & 1 & $\hat{G}_{1}(1)_{10}$ & $\hat{G_{2}}(0,\hat{G}_{1}(0)_{6})_{10}$ & $y_{10}$\\
&11 & 1 & $\hat{G}_{1}(0)_{11}$ & $\hat{G_{2}}(1,\hat{G}_{1}(1)_{3})_{11}$ & $y_{11}$\\
&12 & 1 & $\hat{G}_{1}(0)_{12}$ & $\hat{G_{2}}(1,\hat{G}_{1}(0)_{4})_{12}$ & $y_{12}$\\
&13 & 1 & $\hat{G}_{1}(0)_{13}$ & $\hat{G_{2}}(0,\hat{G}_{1}(1)_{5})_{13}$ & $y_{13}$\\
&14 & 1 & $\hat{G}_{1}(0)_{14}$ & $\hat{G_{2}}(0,\hat{G}_{1}(0)_{6})_{14}$ & $y_{14}$\\
&15 & 0 & $\hat{G}_{1}(1)_{15}$ & $\hat{G_{2}}(1,\hat{G}_{1}(1)_{3})_{15}$ & $y_{15}$\\
&16 & 0 & $\hat{G}_{1}(1)_{16}$ & $\hat{G_{2}}(1,\hat{G}_{1}(0)_{4})_{16}$ & $y_{16}$\\
&17 & 0 & $\hat{G}_{1}(1)_{17}$ & $\hat{G_{2}}(0,\hat{G}_{1}(1)_{5})_{17}$ & $y_{17}$\\
&18 & 0 & $\hat{G}_{1}(1)_{18}$ & $\hat{G_{2}}(0,\hat{G}_{1}(0)_{6})_{18}$ & $y_{18}$\\
&19 & 0 & $\hat{G}_{1}(0)_{19}$ & $\hat{G_{2}}(1,\hat{G}_{1}(1)_{3})_{19}$ & $y_{19}$\\
&20 & 0 & $\hat{G}_{1}(0)_{20}$ & $\hat{G_{2}}(1,\hat{G}_{1}(0)_{4})_{20}$ & $y_{20}$\\
&21 & 0 & $\hat{G}_{1}(0)_{21}$ & $\hat{G_{2}}(0,\hat{G}_{1}(1)_{5})_{21}$ & $y_{21}$\\
&22 & 0 & $\hat{G}_{1}(0)_{22}$ & $\hat{G_{2}}(0,\hat{G}_{1}(0)_{6})_{22}$ & $y_{22}$\\
\addlinespace[3pt]
\noalign{\hrule height 1pt}
\end{tabularx}

\vspace{4pt}
{\footnotesize
\emergencystretch=3em
\setlength{\parindent}{1.5em}
\textsuperscript{1} $\hat{M}_{1}(1)$ and $\hat{M}_{1}(0)$ represent the distributions of insomnia in
arms 1 and 2, respectively. $\hat{G}_{1}(1)$  and $\hat{G}_{1}(0)$ represent the distributions of assigned
insomnia from random draws from $\hat{M}_{1}(1)$ and $\hat{M}_{1}(0)$, respectively. 

$\hat{M}_{2}(1,\hat{G}_{1}(1))$, $\hat{M}_{2}(1,\hat{G}_{1}(0))$,
$\hat{M}_{2}(0,\hat{G}_{1}(1))$, and $\hat{M}_{2}(0,\hat{G}_{0}(1))$ represent the distributions 
of headache in arms 3-6. $\hat{G}_{2}(1,\hat{G}_{1}(1))$, $\hat{G}_{2}(1,\hat{G}_{1}(0))$,
$\hat{G}_{2}(0,\hat{G}_{1}(1))$, and $\hat{G}_{2}(0,\hat{G}_{0}(1))$ represent the distributions of assigned 
headache from random draws from the corresponding empirical distributions of $M_{2}$. 

Finally, $y$ gives the observed average level of depression in each arm.\par}
\vspace{10pt}
% -------

Such an experiment fulfills all the requirements of Assumption~\ref{assm:seq_ignorability_2med}, so any
estimand written in terms of $\Psi(a,a',a'',a''')$ can be identified. 

\subsection{Comparison with Previous Definitions of Interventional Effects for Multiple Mediators}

We now show how our definition of interventional effects for two causally ordered mediators compares to 
a number of different ways that interventional effects for multiple mediator problems have been defined in 
the literature. 

\subsubsection{~\cite{linvanderweele2017}}

~\cite{linvanderweele2017} provide the most similar definition of interventional effects to ours, and 
equation~\ref{eq:RIE_two_mediator_g_comp} and Assumption~\ref{assm:seq_ignorability_2med} are adapted
from this paper. A similar type of path-specific effect is also discussed 
in~\citep{vanderweeletchetgentchetgen2017}. The main idea in~\cite{linvanderweele2017} is to estimate
path-specific interventional effects in the multiple mediator setting in a similar way to how 
such effects are estimated in the single mediator case. In particular, they pursue a decomposition
of the total interventional effect in terms of path-specific interventional effects 
$A \rightarrow Y$, $A \rightarrow M_{1} \rightarrow Y$, $A \rightarrow M_{2} \rightarrow Y$, 
and $A \rightarrow M_{1} \rightarrow M_{2} \rightarrow Y$, which mirrors the decomposition of 
the total interventional effect in terms of the pure direct interventional effect and the total 
indirect interventional effect in the one mediator case. This decomposition relies on the following
simple algebraic expansion: 

\begin{equation}
\begin{aligned}
	\mathbb{E}[Y_{i}(1,G_{1,i}(1),G_{2,i}(1,G_{1,i}(1)))] - 
	\mathbb{E}[Y_{i}(0,G_{1,i}(0),G_{2,i}(0,G_{1,i}(0)))]\\ 
	= \Psi(1,1,1,1) - \Psi(0,0,0,0)\\
	= \{\Psi(1,0,0,0) - \Psi(1,0,0,0)\} + \{\Psi(1,1,0,0) - \Psi(1,0,0,0)\} +\\ 
	\{\Psi(1,1,1,0) - \Psi(1,1,0,0)\} + \{\Psi(1,1,1,1) - \Psi(1,1,1,0)\}
\end{aligned}
\end{equation}

Each term represents a particular interventional effect that corresponds to the pathways 
$A \rightarrow Y$, $A \rightarrow M_{1} \rightarrow Y$, $A \rightarrow M_{2} \rightarrow Y$, 
and $A \rightarrow M_{1} \rightarrow M_{2} \rightarrow Y$, respectively. 

\subsubsection{~\cite{vansteelandtdaniel2017}}

\subsubsection{~\cite{vanderweelerobinson2014} and~\cite{jacksonvanderweele2018}}

\subsubsection{~\cite{vanderweeleetal2014} for Effects Through $M_{2}$}

\section{Interpretation of Interventional Effects}\label{subsection:interpretation}

\subsection{Experimental Manipulability}

% Pure vs. total indirect comes from Robins and Greenland 1992
% Robins and Greenland 1992 interaction example (can drop to binary mediator/outcome if helpful)
% Potential table of attribution for each effect
% Vanderweele decomposition


% Emphasize that these are rough definitions
% Main sources for the development of these effects: Robins and Greenland (1992); Pearl (2001)
% Imai, Keele, Yamamoto 2010: sequential ignorability assumption (read this paper again closely)
% Imai, Tingley, Yamamoto 2013 on experimental design? Sensitivity analysis?
% Various citations for issues with natural effects: Vansteelandt and Vanderweele 2009 (interpretation/un-testability), 
% Vanderweele et al 2014 and Avin et al (exposure-induced confounding) 

% Section 4: look at Robins and Greenland 1992, do a similar breakdown of the types of patients in a binary setting 
% in the two ordered mediators case; generalize, using Vanderweele 2013/2014, to the interventional effects case

% Should the Vanderweele decomposition framework enter the picture in Section 3? This would allow me to use the framework
% when discussing natural effects from the beginning, but it also would require a serious re-write as well as an 
% acknowledgement that some of the effects definitions (which ones?) don't make much sense (except that they do!)

% I should do this, but how to frame?

% Structure:
% - Potential outcomes framework
% - Mediation analysis with one mediator 
%.     - Notation, causal story (informal direct and indirect effects) (Pearl 2001)
%.     - Natural effects (pure vs. indirect )
%             - Breakdown of types of patients in the applied example (Robins and Greenland 1992)
%             - Mapping these types of patients to pure vs. natural direct/indirect effect 
%               definitions (Robins and Greenland 1992; Pearl 2001)
%             - Handling interaction terms (Vanderweele 2013/2014)
%             - Assumptions: standard ones + sequential ignorability, causal graph (Imai et al 2010)
%             - Problems: interpretational (Vanderweele and Vansteelandt 2009; briefly Robins and Greenland 2003) 
%               and mathematical (cross-world independence; exposure induced confounding) 
%               (Imai et al 2010; Avin et al 2005; Vanderweele et al 2014)
%.            - Mention/define controlled direct effects somewhere; this can be part of the 4-way decomposition from
%               Vanderweele 2014
%.            - Emphasize throughout (citing Robins and Greenland, Vanderweele stuff) the importance of the substantive
%               properties of the problem; this will pre-emptively justify the approach taken in Section 4
%.     - Interventional effects with one mediator 
%
% I should incorporate an applied example (probably medical is easiest) throughout my definition of the effects; will make
% it easiest to break down the sub-types of effects like in Robins and Greenland and the Vanderweele papers
% Social mobility example? I don't think this is the best because of the lack of credibility of a causal story. Maybe an
% epi example? Moreno-Betancur headache/anti-depressant causal story seems the best, because it can generalize to Section 4
%
% Differences between Section 3 and Section 4: Section 3 is a pure interpretation of the existing literature while
% Section 4 attempts to advance a "novel" framework that can clarify/make commensurable the various 
% multiple mediator effect definitions 

% Notes/To-do:
% - Additive vs. proportional scale (throat-clearing)
% - Composition assumption (Vanderweele and Vansteelandt, 2009; Pearl, 2001)
% - Need to check binary vs. continuous mediator and outcome; consider at the very least switching to binary mediator, then
% giving general equations for MC estimation in section 5
% - Wording: treatment vs. exposure

% Independence of Y(a,m) and M